% !PRINTED 213
\weberpage{213}
%<<CONT
are, and this equation cannot split even when the domain of rationality of the rational numbers is extended by adjoining the independent variable $\varkappa^2$. % [source syntax] German prints the sentence with malformed word order.

The degree of this equation is $\varphi(n)$, i.\,e. equal to the number of the numbers $m$ incongruent modulo $n$ that are relatively prime to $n$.

Now let

\weq{13}{F'(r) = 0} % [sic] The German source prints F' here; the subsequent argument consistently uses F.

be a \gesperrt{Galois} resolvent of degree $\mu$ of the division equation in the original domain of rationality (so that all roots $x_{\nu,\nu'}$ can be represented rationally by $r$, Bd.\,I, § 152). This equation must \gesperrt{split} after the adjunction of a root of (12); for if it did not, then every rational relation

\weq{14}{\psi(r,\varrho) = 0}

would be satisfied for \gesperrt{all} roots of (13), and $\psi(t,\varrho)$ would be divisible by $F(t)$ (for a variable $t$). Hence (14) would still continue to hold if $\varrho$ were replaced by another root $\varrho^m$ of (12). But by § 62, (12) this is not the case when one puts one of the \gesperrt{Abelian} relations in place of (14).

Now let, \gesperrt{after the adjunction of} $\varrho$,

\weq{15}{F(r,\varrho) = 0}

be the \gesperrt{Galois} resolvent of degree $\nu$ of the division equation, so that, by 2., $\nu$ is equal to the order of the group $\mathfrak{B}$. Then $F(t)$ is algebraically divisible by $F(t,\varrho)$, and hence, because of the irreducibility of (12), also by each of the functions $F(t,\varrho^m)$. Since the functions $F(t,\varrho^m)$ are irreducible, two of them can have a common divisor only when they are entirely identical. But if

\[
F(t,\varrho) = F(t,\varrho^m)
\]

were the case, then from every equation of the form (14) it would follow that $\psi(t,\varrho)$ is divisible by $F(t,\varrho) = F(t,\varrho^m)$, and it would follow that in (14) the \gesperrt{interchange} $(\varrho,\varrho^m)$ is admitted, which again is not the case for the \gesperrt{Abelian} relations. Consequently the $\varphi(n)$ functions $F(t,\varrho^m)$ are all different from one another, and $F(t)$ is divisible by their product. The degree $\mu$ of $F(t)$ is therefore at least $= \nu\varphi(n)$. But it also cannot be higher than $\nu\varphi(n)$, since $\nu\varphi(n)$ is the order of the group $\mathfrak{A}$, and the
%CONT>>
